\scanpage{88}
Les isom de transitivité se \unsure{décomposent}
en trois isomorphismes de transitivité
particuliers :

a) deux morphismes lisses

b) Deux morphismes d'immersion,

c) commutativité des images inverses
dans le carré ($\unsure{Q|2}$). Les cas a) b)
triviaux, le cas du carré $\unsure{Q|2}$ est essentiellement
le th. de changement de base lisse.

\uline{NB.}\quad Dans tous ces isom de
transitivité, il faudrait vérifier que
quand il y a trois composés, on a les
\unsure{compatibilités|compatibilité} habituelles --- Nous l'admettrons.

\deleted{\unsure{Isom}}
2$^\circ$)\quad \uline{Formules d'adjonction}
\diagram{\begin{tikzcd}[ampersand replacement=\&,column sep=small]
X\arrow[d,no head,"f"]\&F^\bullet\\
Y\&G^\bullet
\end{tikzcd}}
\[
\boxed{\operatorname{Hom}(\mathbb{R}_!f(F^\bullet),G^\bullet)
 \simeq\operatorname{Hom}(F^\bullet,\mathbb{R}^!f(G^\bullet))}
\]
\[
[F^\bullet\in\operatorname{Ob}\mathrm D^b(X),\quad
 G^\bullet\in\operatorname{Ob}\mathrm D^+(Y)]
\]
\[
\boxed{\operatorname{Hom}(\mathbb{L}f^*(G^\bullet),F^\bullet)
 \simeq\operatorname{Hom}(G^\bullet,\mathbb{R}f_*(F^\bullet))}
\]
\[
[F^\bullet\in\operatorname{Ob}\mathrm D^+(X),\quad
 G^\bullet\in\operatorname{Ob}\mathrm D^b(Y)]
\]
\mspage{87}
